\documentclass[12pt]{article}
\usepackage{fullpage}
\usepackage{amsmath,amssymb,amsthm}

\newtheorem{lemma}{Lemma}
\newtheorem{proposition}{Proposition}
\newtheorem{theorem}{Theorem}

\newcommand{\Dil}{\operatorname{Dil}}
\newcommand{\Vol}{\operatorname{Vol}}
\newcommand{\Fill}{\operatorname{Fill}}
\newcommand{\Mass}{\operatorname{Mass}}
\newcommand{\RR}{\mathbb R}

\begin{document}

\begin{center}
{\bf A reduction for the rectangle \(2\)-dilation estimate}
\end{center}

\section*{Problem statement and interpretation}

Rectangles are oriented Euclidean boxes.  A map
\(f:(R,\partial R)\to(S,\partial S)\) has degree \(1\) if
\(f_*[R,\partial R]=[S,\partial S]\) in relative homology.  I denote the
small constant in the problem by \(\kappa\), reserving \(k\) for dilation
dimension.  Constants below are positive dimensional constants.  Strict
inequalities in the problem follow from non-strict inequalities after decreasing
\(\kappa\).

This manuscript proves the standard first alternative, the sharp reductions, and a
self-absorbing plaque-puncture estimate which is the right replacement for the
false single-puncture lemma from earlier drafts.  One parametric tightening theorem
remains unproved here; it is isolated explicitly in the last section.

\section*{Standard tools}

We use integral Lipschitz currents.  Let \(T=[R,\partial R]\).  After an arbitrarily
small PL/Lipschitz approximation, fixed on a collar of \(\partial R\) and preserving
the relative degree, the target coordinate projections used below may be assumed
transverse to the relevant planes; the estimates then pass to the original map by
lower semicontinuity of mass and dilation.  For a projection \(\pi\) and
\(F=\pi\circ f\), slicing naturality gives, for a.e. regular value \(y\),
\[
        f_\#\langle T,F,y\rangle
        =\langle [S,\partial S],\pi,y\rangle
\]
as relative currents.

If \(\Dil_2(f)\le1\), then \(\Dil_j(f)\le1\) for \(j=3,4\): if
\(\sigma_1\ge\cdots\ge\sigma_4\) are the singular values of \(df\), then
\(\sigma_1\sigma_2\le1\), hence \(\sigma_i\le1\) for \(i\ge2\), and
\(\sigma_1\cdots\sigma_j\le1\) for \(j\ge2\).

We use two estimates of Guth.  First, the \(k=2,n=4\) relative isoperimetric
profile theorem for rectangles says that if \(Z\) is a relative integral \(2\)-cycle
in \(R\), \(A=\Mass Z\le c_I R_1R_2\), then
\[
        \Fill_R(Z)\le C_I
        \max\left\{ A^{3/2},\,\frac{A^2}{R_1}\right\}.        \tag{1}
\]
Second, Guth's degree estimate with \(j=1,k=2,l=4\) gives
\[
        R_1^3R_2R_3R_4\ge c_G S_1^3S_2S_3S_4                  \tag{2}
\]
for every degree-one map of pairs with \(\Dil_2(f)\le1\).

\section*{The first alternative}

\begin{lemma}[central target filling]\label{lem:target-fill}
Let
\[
 P_y=[0,S_1]\times[0,S_2]\times\{y_3\}\times\{y_4\},
\]
where \(y_3\in[S_3/3,2S_3/3]\) and
\(y_4\in[S_4/3,2S_4/3]\).  Every relative \(3\)-chain \(Y\) in \(S\) with
\(\partial Y=P_y\) in \(I_2(S,\partial S)\) obeys
\[
        \Mass Y\ge c\,S_1S_2S_3 .
\]
\end{lemma}

\begin{proof}
Choose a smooth \(1+o(1)\)-Lipschitz function \(\psi\) on the
\((x_3,x_4)\)-rectangle, zero on its boundary and satisfying
\(\psi(y_3,y_4)\ge cS_3-o(1)\).  Apply Stokes for relative currents to
\(\omega=dx_1\wedge dx_2\wedge d\psi\).  The comass of \(\omega\) is at most
\(1+o(1)\), and
\[
 \int_Y\omega=\int_{\partial Y}\psi\,dx_1\wedge dx_2 .
\]
The part of \(\partial Y\) lying in \(\partial S\) contributes zero: on the
\(x_1,x_2\)-faces \(dx_1\wedge dx_2=0\), while on the \(x_3,x_4\)-faces
\(\psi=0\).  Letting \(o(1)\to0\) proves the claim.
\end{proof}

\begin{proposition}[first alternative]\label{prop:first}
There are constants \(c_1,C_1>0\) such that the following holds.  Let
\(f:(R,\partial R)\to(S,\partial S)\) have degree \(1\) and \(\Dil_2(f)\le1\).
If
\[
        C_1R_1R_2^2\le S_1S_2S_3 ,
\]
then \(R_3R_4\ge c_1S_3S_4\).
\end{proposition}

\begin{proof}
Let \(F=(f_3,f_4)\) and
\(Q=[S_3/3,2S_3/3]\times[S_4/3,2S_4/3]\).  For a.e. \(y\in Q\), put
\(Z_y=\langle [R,\partial R],F,y\rangle\).  Slicing naturality gives
\(f_\#Z_y=P_y\).  Any relative filling of \(Z_y\) pushes forward to a relative
filling of \(P_y\), and \(\Dil_3(f)\le1\).  Lemma \ref{lem:target-fill} gives
\[
        \Fill_R(Z_y)\ge cS_1S_2S_3 .                         \tag{3}
\]
If \(C_1\) is large enough, then (1) and (3) force
\(\Mass Z_y\ge cR_1R_2\) for a.e. \(y\in Q\).  Coarea and \(J_2F\le1\) give
\[
 \Vol(R)\ge\int_Q\Mass Z_y\,dy
      \ge c R_1R_2S_3S_4 ,
\]
and hence \(R_3R_4\ge cS_3S_4\).
\end{proof}

\section*{Consequences of the same slicing}

The preceding argument also gives, for a.e. \(y\in Q\),
\[
 \Mass Z_y\ge c\min\left\{R_1R_2,\,
          (S_1S_2S_3)^{2/3},\,
          (R_1S_1S_2S_3)^{1/2}\right\}.                       \tag{4}
\]
Indeed, if \(A=\Mass Z_y<c_I R_1R_2\), then (1) and (3) imply
\[
 cS_1S_2S_3\le C_I\max\{A^{3/2},A^2/R_1\}.
\]
After integration over \(Q\),
\[
 \Vol(R)\ge cS_3S_4\min\left\{R_1R_2,\,
          (S_1S_2S_3)^{2/3},\,
          (R_1S_1S_2S_3)^{1/2}\right\}.                      \tag{5}
\]
Also, (2) and \(R_1\le\kappa S_1\) imply
\[
        \Vol(R)\ge c_G\kappa^{-2}S_1S_2S_3S_4 .              \tag{6}
\]
Thus (6) proves the desired second alternative when
\(S_3/S_2\le c\,\kappa^{-6}\).

Assume now that the first desired alternative fails and take
\(\kappa<\min\{c_1,1\}\).  Proposition \ref{prop:first} gives
\[
        R_1R_2^2\gtrsim S_1S_2S_3 .                          \tag{7}
\]
Writing \(\alpha=R_1/S_1\le1\) and
\(T=S_1S_2^{1/2}S_3^{3/2}S_4\), (5) gives
\[
        \Vol(R)\ge c\,\alpha^{1/2}T,                          \tag{8}
\]
while (2) gives
\[
        \Vol(R)\ge c\,\alpha^{-2}S_1S_2S_3S_4 .                \tag{9}
\]
Consequently, if the second desired alternative also fails, then necessarily
\[
        \alpha\lesssim \kappa^2,\qquad
        \frac{S_3}{S_2}\gtrsim \alpha^{-4}\kappa^{-2}.          \tag{10}
\]
Thus any remaining counterexample must be extremely anisotropic.

\section*{Weighted coarea}

Let \(G=(f_1,f_2)\).  At rank-two points of \(F=(f_3,f_4)\), set
\[
        \lambda(x)=\left\|dG|_{\ker dF_x}\right\|,
\]
and set \(\lambda=0\) on the rank-deficient set.  If
\(E_y=\int_{Z_y}\lambda^2\,d\|Z_y\|\) and \(I=\int_QE_y\,dy\), then
\[
        I\le \Vol(R),                                      \tag{11}
\]
and, for a.e. \(y\in Q\),
\[
        E_y\ge S_1S_2 .                                    \tag{12}
\]
For (11), write \(T_x=\ker dF_x\), \(N_x=T_x^\perp\).  If
\(\lambda>0\), choose a unit \(v\in T_x\) with \(|dG(v)|=\lambda\).  For every unit
\(n\in N_x\),
\[
        |df(v)\wedge df(n)|\ge \lambda |dF(n)|.
\]
Since \(\Dil_2(f)\le1\), the two singular values of \(dF|_{N_x}\) are at most
\(\lambda^{-1}\), so \(\lambda^2J_2F\le1\).  Coarea gives (11).  For (12),
slicing naturality gives \(G_\#Z_y=[0,S_1]\times[0,S_2]\) as a relative current; the
area formula and \(J_2(G|_{Z_y})\le\lambda^2\) give the claim.

\section*{The self-absorbing plaque-puncture reduction}

Let
\[
        A=\int_Q\Fill_R(Z_y)\,dy,\qquad q=|Q|\sim S_3S_4 .
\]
Choose measurable fillings \(B_y\) with \(\partial B_y=Z_y\) in \(I_2(R,\partial R)\)
and \(\int_Q\Mass B_y\,dy\le2A\).  For \(i<j\), let \(P_{ij}(x)\) be the coordinate
\(2\)-plaque through \(x\), varying only \(x_i,x_j\).  For \(z\in Q\), define
\[
        \mathcal B_z=\{x\in R:\ z\in F(P_{ij}(x))\hbox{ for some }i<j\}.
\]

\begin{lemma}[all-plaque bad mass]\label{lem:plaque}
\[
 \frac1q\int_Q\int_Q \Mass(B_y\llcorner\mathcal B_z)\,dy\,dz
 \le C\,\frac{R_3R_4}{S_3S_4}\,A .
\]
\end{lemma}

\begin{proof}
For fixed \(x\), the set of \(z\) with \(x\in\mathcal B_z\) is contained in
\(\bigcup_{i<j}F(P_{ij}(x))\cap Q\).  Since \(J_2(F|_{P_{ij}})\le1\) and
\(R_iR_j\le R_3R_4\), this set has area at most \(C R_3R_4\).  Integrating the
indicator of \(\mathcal B_z\) with respect to \(d\|B_y\|(x)\,dy\,dz\) and using
\(\int_Q\Mass B_y\,dy\le2A\) proves the estimate.
\end{proof}

The estimate needed to finish the theorem is now the following fixed-puncture
parametric tightening statement:
\[
 A\le C\left(R_1I+\frac{I^2}{S_1q}\right)
      +C\int_Q\Mass(B_y\llcorner\mathcal B_z)\,dy              \tag{13}
\]
for at least one \(z\in Q\), or equivalently after averaging for a random \(z\).
Together with Lemma \ref{lem:plaque}, (13) gives
\[
 A\le C\left(R_1I+\frac{I^2}{S_1q}\right)
      +C\frac{R_3R_4}{S_3S_4}A.                              \tag{14}
\]
In the hard regime \(R_3R_4\le\kappa S_3S_4\), choosing \(\kappa\) small absorbs the
last term.  Since (3) gives \(A\ge cS_1S_2S_3q\), either
\(R_1I\gtrsim S_1S_2S_3q\), which with \(R_1\le\kappa S_1\) implies
\(I\gtrsim T\), or
\[
        \frac{I^2}{S_1q}\gtrsim S_1S_2S_3q,
\]
which gives \(I\gtrsim T\) directly.  Then (11) gives
\(\Vol(R)\ge cT\), the desired second alternative.

\section*{Directional filling check}

For a rectifiable relative \(2\)-cycle \(Z\subset R\), define
\[
        \Vol_{ij}(Z)=\int_Z |dx_i\wedge dx_j|\,d\|Z\|.
\]
Proposition 2.2 of Guth's directional-isoperimetric paper, specialized to
\(m=2,n=4\), gives a relative filling \(Y\) of \(Z\) with
\[
 \Mass Y\le C\{R_1(\Vol_{23}Z+\Vol_{24}Z+\Vol_{34}Z)
      +R_2(\Vol_{13}Z+\Vol_{14}Z)+R_3\Vol_{12}Z\}.            \tag{15}
\]
Indeed, the four triples give respectively
\[
\begin{array}{c|c}
123& R_1\Vol_{23}+R_2\Vol_{13}+R_3\Vol_{12}\\
124& R_1\Vol_{24}+R_2\Vol_{14}\\
134& R_1\Vol_{34}\\
234&0 .
\end{array}
\]

Applying (15) to \(Z_y\) and integrating, (3) gives
\[
 S_1S_2S_3\,q
 \le C\{R_1(M_{12}+M_{13}+M_{14})
      +R_2(M_{23}+M_{24})+R_3M_{34}\},                       \tag{16}
\]
where
\[
        M_{ij}=\int_{F^{-1}(Q)}
        \left|\frac{\partial(f_3,f_4)}{\partial(x_i,x_j)}\right|\,dx .
\]
The bookkeeping is complementary: for example,
\(\int_Q\Vol_{12}(Z_y)\,dy=M_{34}\), because
\(dx_1\wedge dx_2\wedge df_3\wedge df_4\) contains the \(x_3x_4\)-minor of \(F\).
The absolute-value minor inequalities (16), even together with the known monomial
degree estimates, do not imply (13); the latter is genuinely parametric.

\section*{Remaining open issues}

The proof is incomplete at exactly one point: the fixed-puncture parametric
tightening estimate (13).  The all-plaque bad set in Lemma \ref{lem:plaque} fixes
the earlier false random-puncture lemma: linear examples in which \(F\) varies only
in the \(x_1,x_2\)-directions are now charged because the \(x_1x_2\)-plaques detect
the puncture.  What remains is to turn the fact that \(\rho_z\circ F\) is defined on
the good coordinate \(2\)-skeleton into controlled fillings whose good cost is
\(R_1I+I^2/(S_1q)\).  A proof would likely be a cubical Federer--Fleming tightening
over the parameter square \(Q\), with the \(R_1I\) term coming from pushing in the
short source direction and the quadratic term from a coarea optimization in the
\(f_1\)-direction.  The present document proves all reductions needed once (13) is
established, but does not establish (13).

\begin{thebibliography}{9}
\bibitem{GuthExpanding}
L. Guth, \emph{Area-expanding embeddings of rectangles}, arXiv:0710.0403.

\bibitem{GuthDirectional}
L. Guth, \emph{Directional isoperimetric inequalities and rational homotopy
invariants}, arXiv:0802.3549.

\bibitem{GuthThesis}
L. Guth, \emph{Area-contracting maps between rectangles}, Ph.D. thesis, MIT, 2005.

\bibitem{Federer}
H. Federer, \emph{Geometric Measure Theory}, Springer, 1969.
\end{thebibliography}

\end{document}
